%% Chen-Yu Kuo — resume (English)  [v4: single-line date ranges; rebalanced away from EEW]
%% Build: tectonic -X compile cv-en.tex
%%
%% Layout uses resume.cls by Weitian LI (LPPL 1.3c), vendored here verbatim:
%%   https://github.com/liweitianux/resume
%% itself based on YACC by Christophe Roger and Plasmati Graduate CV.

\documentclass{resume}

%% Density tuning. resume.cls defaults to margin=1.5cm and \linespread{1.15};
%% this CV carries more content, so tighten a little without changing the look.
\usepackage{geometry}
\geometry{top=1.0cm, bottom=0.9cm, left=1.05cm, right=1.05cm, includefoot}
\linespread{1.0}
\setlength{\parskip}{0.35ex plus 0.3ex minus 0.1ex}
\setlist[itemize,1]{label=\faAngleRight, leftmargin=1.2em, labelsep=0.45em,
                    topsep=0.2ex, itemsep=0.25ex, parsep=0pt}
\titlespacing{\section}{0em}{0.6ex}{0.6ex}

%% v4: dates are written as one-line ranges ("Nov 2023 -- Nov 2024") instead of the
%% class's stacked end-over-start pair, which read as two separate dates. The widest
%% range measures 96.4pt, so the date column needs 81pt -> 103pt; the right column
%% gives back the same 22pt so the total width is unchanged.
\addtolength{\leftcolwidth}{22pt}
\addtolength{\rightcolwidth}{-22pt}

%% Scale the class's fonts down a touch so the content lands on two pages.
\setmainfont{IBM Plex Serif}[Scale=0.93]
\setmonofont{IBM Plex Mono}[Scale=0.87]

%% resume.cls only loads xeCJK under the [zh] option; this English document
%% still needs CJK glyphs for a handful of proper nouns.
\usepackage{xeCJK}
\setCJKmainfont{Noto Serif CJK TC}[Scale=0.93]
\XeTeXlinebreaklocale "zh"

\fileinfo{%
  \faGithub{} \link{https://github.com/whes1015/cv}{whes1015/cv} \hspace{0.5em}
  Layout: \link{https://github.com/liweitianux/resume}{liweitianux/resume} (LPPL~1.3c) \hspace{0.5em}
  \faEdit{} \today
}

\name{Chen-Yu}{Kuo}
\keywords{Go, TypeScript, Python, C++, Rust, Flutter, eBPF, XDP, CoreDNS, embedded, ESP32,
  PyTorch, LLM quantisation, seismology, earthquake early warning}
\tagline{郭宸毓 \enspace\textbullet\enspace Systems \& full-stack engineer --- eBPF/Go infrastructure,
  ESP32 firmware, applied ML and LLM inference, earthquake early warning}

\profile{
  \email{whes1015@gmail.com}
  \github{whes1015}
  \iconlink{\faRobot}{https://huggingface.co/YuYu1015}{\texttt{huggingface.co/YuYu1015}} \\
  \degree{B.S. in Electrical Engineering}
  \university{National Kaohsiung Univ. of Science and Technology}
  \address{Kaohsiung, Taiwan}
}

\begin{document}
\makeheader

%======================================================================
\begin{abstract}
I build systems end to end, usually as the only author: an eBPF/XDP packet filter and a five-city
Nginx/CoreDNS fleet written in Go, ESP32 firmware written register by register, Flutter and
TypeScript clients, and PyTorch models I train, evaluate and quantise myself. Most of it runs under
\textbf{ExpTech} (探索科技), the 29-person volunteer organisation I founded and one of two
administrators of, which operates Taiwan's largest independent earthquake early-warning platform.
Alongside it I worked as a backend engineer at P-Waver, then as a part-time assistant at Academia
Sinica's Institute of Information Science, and now write software at the National Health Research
Institutes.
\end{abstract}

%======================================================================
\sectionTitle{Competences}{\faWrench}
%======================================================================
\begin{competences}[9em]
  \comptence{Languages}{%
    TypeScript/JavaScript, Go, Python, Dart, C/C++, Rust, Lua, Shell
  }
  \comptence{Infrastructure}{%
    Linux, eBPF/XDP, Docker, Nginx/OpenResty, CoreDNS, NTP/SNTP, Prometheus, CI/CD
  }
  \comptence{Backend \& apps}{%
    Go and Node/TypeScript services, REST/WebSocket/SSE, Flutter (iOS \& Android), Electron, MapLibre
  }
  \comptence{Embedded}{%
    ESP32/Arduino, FreeRTOS, I\textsuperscript{2}C/SPI, LSM6DS3, GNSS (QZSS L1S), LoRa/Meshtastic
  }
  \comptence{Machine Learning}{%
    PyTorch, SeisBench, ONNX, scikit-learn, vLLM, llm-compressor, TensorRT Model Optimizer
  }
  \comptence{Domain}{%
    Earthquake early warning, GMPE calibration, CWA/JMA intensity scales, phase picking
  }
  \comptence{\icon{\faLanguage} Languages}{%
    \textbf{Mandarin} --- native; \textbf{English} --- technical reading \& writing
  }
\end{competences}

%======================================================================
\sectionTitle{Education}{\faGraduationCap}
%======================================================================
\begin{educations}
  %% TODO: replace "present" with your expected graduation, e.g. "2022 -- Jun 2027"
  \education%
    {}[2022 -- present]%
    {National Kaohsiung Univ. of Science \& Technology}%
    {Dept. of Electrical Engineering}%
    {Electrical Engineering}%
    {B.S.}
\end{educations}

%======================================================================
\sectionTitle{Experience}{\faBriefcase}
%======================================================================
\begin{experiences}
  %% The whole range goes in the end-date slot and the start-date slot is left off:
  %% resume.cls only closes the table row when a description or tag list follows, so
  %% \experience[start]{end}{summary} alone would merge into the next row.
  \experience%
    {Jan 2025 -- present}%
    {\textbf{Software Developer} --- National Health Research Institutes
      (國家衛生研究院), Taiwan}

  \separator{0.6em}

  \experience%
    {Dec 2024 -- Jun 2025}%
    {\textbf{Part-time Assistant} --- Institute of Information Science,
      Academia Sinica (中央研究院資訊科學研究所), Taipei}

  \separator{0.6em}

  \experience%
    {Nov 2023 -- Nov 2024}%
    {\textbf{Backend Engineer} --- P-Waver Inc.
      (衛波科技, earthquake early-warning systems), Taipei}

  \separator{0.6em}

  \experience%
    {Oct 2020 -- present}%
    {\textbf{Founder \& Lead Engineer} --- ExpTech (探索科技), Taiwan}%
    [\begin{itemize}
      \item Lead a 29-member volunteer organisation across 234 repositories as one of its two
        administrators, spanning firmware, 37 Go services, Flutter clients and ML pipelines.
        Authored \textbf{13,411} commits inside it (19,284 across GitHub, 477 PRs) and am top
        contributor on every flagship codebase --- DPIP 838/2,339, TREM-Lite 448/847, platform API
        480/803 --- and sole author of the 665-commit production infrastructure repository.
      \item Signed a data agreement with Taiwan's \textbf{Central Weather Administration (CWA)} in
        late 2022 --- reported by \emph{親子天下} as the Seismological Center's first non-corporate
        partner; invited back 10+ times as a collaborator (\emph{天下雜誌}, Feb 2026).
      \item Ship under real stakes: \textbf{100,000+} Google Play installs and 4.2/5 from 1,378
        App Store ratings for DPIP; \textbf{328,542} installer downloads across TREM desktop
        releases. During the M7.2 Hualien earthquake (3 Apr 2024) the network recorded 800+ events
        in a single day.
    \end{itemize}]
\end{experiences}

%======================================================================
\sectionTitle{Security, Open Source \& Press}{\faLock}
%======================================================================
\begin{itemize}
  \item \textbf{Security \& incident response} --- operated a honeypot that captured in-the-wild
    exploitation of \textbf{CVE-2025-55182} (React Server Components, CVSS 10.0) within days of its
    December 2025 disclosure; documented six malware samples with verified hashes and fifteen
    attacker source addresses.
  \item \textbf{Open source} --- 12 merged PRs to outside projects: a rendering-performance fix to
    the seismograph suite \textbf{AnyShake}
    (\link{https://github.com/anyshake/spectrogram-js}{\texttt{spectrogram-js}}), a
    Leaflet~$\rightarrow$~MapLibre migration for \texttt{smc.peering.tw}, and Flutter upgrades to
    \textbf{NKUST-AP-Flutter}, my university's student app.
  \item \textbf{Press} --- \emph{天下雜誌}, Feb 2026 (quoted as an ExpTech collaborator in a CWA
    Seismological Center profile); \emph{親子天下}, Jul 2024 (DPIP feature,
    {\raise0.06ex\hbox{$\sim$}}500,000 downloads after the 3 Apr 2024 quake).
\end{itemize}

\newpage
%======================================================================
\sectionTitle{Selected Engineering Work}{\faCode}
%======================================================================
\begin{projects}
  \project%
    {}[Dec 2025 -- present]%
    {sole author}%
    {Go, eBPF/XDP, OpenResty, CoreDNS}%
    {Production infrastructure --- packet filtering, edge and DNS}%
    {\begin{itemize}
      \item \textbf{kekkai} --- packet filter in Go using \texttt{cilium/ebpf}, with an XDP ingress
        program and a TC egress hook; \textsc{lru\_hash}/\textsc{array}/\textsc{ringbuf} maps drive
        per-IP rate limiting, allow/blocklists and port policy, with an emergency detach path.
      \item \textbf{Nginx/OpenResty fleet} --- 23 upstreams across nodes in Taipei, Kaohsiung,
        Tainan, Tokyo and Osaka, with Lua modules for traffic accounting, NTP and image proxying.
      \item \textbf{Split-horizon CoreDNS} with a custom Go \texttt{dnsstats} plugin
        (\texttt{miekg/dns}) exposing per-domain/client query stats and Prometheus metrics. Sole
        author too of the NTP server, API gateway and Meshtastic LoRa ingest service.
    \end{itemize}}

  \separator{0.7em}

  \project%
    {}[Apr 2026 -- Jul 2026]%
    {sole author}%
    {Python, PyTorch, vLLM, TensorRT Model Optimizer}%
    {LLM inference \& model compression}%
    {\begin{itemize}
      \item Published \textbf{27 quantised models} on Hugging Face
        ({\raise0.06ex\hbox{$\sim$}}\textbf{73,000} all-time downloads) for vLLM and llama.cpp:
        \textbf{NVFP4} via TensorRT Model Optimizer, INT4-AutoRound and GGUF, up to 35B
        mixture-of-experts.
      \item Implemented \textbf{refusal-direction ablation} --- deriving the refusal direction from
        paired prompts and orthogonalising it out of the weight matrices, rather than steering at
        inference time --- and released the \emph{Ornith-1.0} 9B/35B line including a DPO-tuned
        variant.
    \end{itemize}}

  \separator{0.7em}

  \project%
    {}[Sep 2025 -- Jun 2026]%
    {sole author, 7,759 LOC}%
    {C++, ESP32, FreeRTOS}%
    {ES-Net --- sensor firmware \& 200-node device network}%
    {\begin{itemize}
      \item Hand-wrote an \textbf{LSM6DS3 accelerometer driver at register level} instead of using a
        vendor library --- own register map, \textsc{who\_am\_i}, \textsc{ctrl1\_xl},
        \textsc{int1\_ctrl}, FIFO --- plus an IIR filter bank and a binary indexed tree computing a
        running amplitude percentile on-device.
      \item Timing is the hard problem here: SNTP in \textsc{smooth} mode so the clock slews rather
        than steps mid-waveform, plus HW/SW watchdogs, OTA update, and a documented custom binary
        wire protocol with TOTP device authentication.
      \item Runs on \textbf{200+ nodes} nationwide; first-generation units were self-assembled for
        under NT\$800 each.
    \end{itemize}}

  \separator{0.7em}

  \project%
    {}[Aug 2023 -- present]%
    {top contributor, 838/2,339 commits}%
    {Dart/Flutter, TypeScript}%
    {DPIP --- cross-platform mobile app \& backend}%
    {\begin{itemize}
      \item iOS and Android client with push alerting, MapLibre layers and 10 interface languages;
        \textbf{100,000+} Google Play installs, 4.2/5 from 1,378 App Store ratings. Wrote the test
        suite and custom commit-gate tooling outright, and authored the TypeScript backend
        \texttt{dpip-server-ts} in full (122/122 commits).
    \end{itemize}}

  \separator{0.7em}

  \project%
    {}[Aug 2026 -- present]%
    {sole author, {\raise0.06ex\hbox{$\sim$}}12k LOC with unit tests}%
    {C++}%
    {QZSS L1S satellite message decoder}%
    {\begin{itemize}
      \item Decodes QZSS L1S 250-bit SBAS frames end to end --- preamble, 6-bit message type,
        212-bit payload and CRC-24Q --- including type 43 \textbf{DC Report} (JMA disaster
        broadcasts, IS-QZSS-DCR) and type 44 DCX.
    \end{itemize}}

  \separator{0.7em}

  \project%
    {}[Mar 2026 -- Jul 2026]%
    {sole author}%
    {Python, PyTorch, SeisBench, ONNX}%
    {Applied ML --- regression \& sequence models on sensor data}%
    {\begin{itemize}
      \item Refit a regression model on \textbf{45,398} CWA PGA/PGV/intensity observations
        (2000--2025), replacing an equation that over-amplified with magnitude and a site-response
        term that was double-counted, inflating estimates {\raise0.06ex\hbox{$\sim$}}2$\times$. Held-out:
        \textbf{62.0\% exact vs 39.7\%} for the model then in production, MAE
        \textbf{0.400 vs 0.671}, 0\% false alarms.
      \item Trained a 50\,Hz PhaseNet sequence model over a 6,457-line pipeline (ingest, noise
        injection, evaluation, ONNX export). Diagnosed a windowing bug --- the source window was
        shorter than the model window, so the model learned the event always sits mid-window and
        fired a false pick every 45\,s stride.
    \end{itemize}}

  \separator{0.7em}

  \project%
    {}[Jan 2023 -- present]%
    {top contributor, 448/847 commits}%
    {JavaScript, Electron, Rust}%
    {TREM --- desktop platform \& third-party plugin ecosystem}%
    {\begin{itemize}
      \item Cross-platform Electron client with \textbf{328,542 installer downloads} across 74
        releases, shipped continuously since Jan 2023. Designed the plugin system third parties
        extend it through and maintain its registry (417 commits); wrote the
        \texttt{trem-monitor} health service in Rust (105/150).
    \end{itemize}}
\end{projects}

\end{document}
